\documentclass[11pt]{article}
% ======================================================================
%  Preambule commun - Sujets Bac Serie C (Physique-Chimie)
%  Style sobre : noir et blanc, sans couleurs, sans filigrane,
%  sans en-tetes ni pieds decoratifs. Figures reconstruites en TikZ.
% ======================================================================
\usepackage[utf8]{inputenc}
\usepackage[T1]{fontenc}
\usepackage{lmodern}
\usepackage{textcomp}
\usepackage{amsmath,amssymb}
\usepackage{enumitem}
\usepackage[a4paper,margin=2.3cm]{geometry}
\usepackage{fancyhdr}
\usepackage{tikz}
\usetikzlibrary{arrows.meta,calc,decorations.pathmorphing,patterns,angles,quotes}

% Symbole degre robuste + justification souple
\DeclareUnicodeCharacter{00B0}{\ensuremath{^\circ}}
\tolerance=2000
\emergencystretch=2em
\frenchspacing

% En-tete / pied minimalistes
% Pied de page (provenance) + entete corrige + encadres pedagogiques
% ======================================================================
%  Fragment commun a tous les styles (sujets et corriges).
%  Inclus depuis chaque dossier serie : % ======================================================================
%  Fragment commun a tous les styles (sujets et corriges).
%  Inclus depuis chaque dossier serie : \input{../../commun}
%  (la compilation se fait toujours depuis le dossier de la serie).
%
%  Style sobre conserve : noir et blanc uniquement, filets fins,
%  pas d'aplats ni de couleurs — lisible en photocopie.
% ======================================================================

% --- Compression maximale des PDF ------------------------------------
% Le public consulte souvent sur reseaux mobiles : chaque Ko compte.
\pdfcompresslevel=9
\pdfobjcompresslevel=3
\pdfminorversion=5

% --- Adresse publique du site ---------------------------------------
% Point unique a mettre a jour si le domaine change (puis recompiler
% tous les PDF : voir sources/README.md).
\newcommand{\bacsiteurl}{annabac.pages.dev}

% --- Pied de page : provenance discrete + numero de page ------------
% Les PDF circulent detaches du site (messageries, photocopies) : cette
% ligne permet de retrouver la source et rappelle la gratuite.
\pagestyle{fancy}
\fancyhf{}
\renewcommand{\headrulewidth}{0pt}
\fancyfoot[L]{\scriptsize\itshape Bac 242 --- annales gratuites du bac congolais --- \bacsiteurl}
\fancyfoot[R]{\thepage}

% --- Entete structure des corriges -----------------------------------
% Usage : \entetecorrige{2020}{C}{Mathématiques}
% Les sujets, reproductions de documents officiels, gardent \entetesujet.
\newcommand{\entetecorrige}[3]{%
  \begin{center}
    {\scshape Baccalaur\'eat --- S\'erie #2}\\[0.35em]
    {\Large\bfseries #3 --- Session #1}\\[0.4em]
    {\itshape Corrig\'e d\'etaill\'e}
  \end{center}
  \vspace{0.2em}
  \hrule height 0.8pt
  \vspace{1.5pt}
  \hrule height 0.4pt
  \vspace{1.2em}}

% --- Encadres pedagogiques (corriges) --------------------------------
% Filet vertical gauche, texte legerement reduit, aucun fond : sobre,
% photocopiable, et tolere les sauts de page (package framed).
% Usage, avec parcimonie (2-3 encadres par corrige, la ou ils eclairent
% une demarche) :
%   \begin{methode} ... \end{methode}   -> « Méthode »
%   \begin{rappel}  ... \end{rappel}    -> « Rappel de cours »
%   \begin{piege}   ... \end{piege}     -> « Piège classique »
\usepackage{framed}
\newenvironment{pedagobarre}{%
  \def\FrameCommand{{\vrule width 1pt}\hspace{0.9em}}%
  \MakeFramed{\advance\hsize-\width\FrameRestore}}%
 {\endMakeFramed}
\newenvironment{blocpedago}[1]{%
  \par\medskip
  \begin{pedagobarre}\small
  \noindent{\bfseries\scshape #1.}\hspace{0.5em}\ignorespaces
}{%
  \end{pedagobarre}\medskip}
\newenvironment{methode}{\begin{blocpedago}{M\'ethode}}{\end{blocpedago}}
\newenvironment{rappel}{\begin{blocpedago}{Rappel de cours}}{\end{blocpedago}}
\newenvironment{piege}{\begin{blocpedago}{Pi\`ege classique}}{\end{blocpedago}}

%  (la compilation se fait toujours depuis le dossier de la serie).
%
%  Style sobre conserve : noir et blanc uniquement, filets fins,
%  pas d'aplats ni de couleurs — lisible en photocopie.
% ======================================================================

% --- Compression maximale des PDF ------------------------------------
% Le public consulte souvent sur reseaux mobiles : chaque Ko compte.
\pdfcompresslevel=9
\pdfobjcompresslevel=3
\pdfminorversion=5

% --- Adresse publique du site ---------------------------------------
% Point unique a mettre a jour si le domaine change (puis recompiler
% tous les PDF : voir sources/README.md).
\newcommand{\bacsiteurl}{annabac.pages.dev}

% --- Pied de page : provenance discrete + numero de page ------------
% Les PDF circulent detaches du site (messageries, photocopies) : cette
% ligne permet de retrouver la source et rappelle la gratuite.
\pagestyle{fancy}
\fancyhf{}
\renewcommand{\headrulewidth}{0pt}
\fancyfoot[L]{\scriptsize\itshape Bac 242 --- annales gratuites du bac congolais --- \bacsiteurl}
\fancyfoot[R]{\thepage}

% --- Entete structure des corriges -----------------------------------
% Usage : \entetecorrige{2020}{C}{Mathématiques}
% Les sujets, reproductions de documents officiels, gardent \entetesujet.
\newcommand{\entetecorrige}[3]{%
  \begin{center}
    {\scshape Baccalaur\'eat --- S\'erie #2}\\[0.35em]
    {\Large\bfseries #3 --- Session #1}\\[0.4em]
    {\itshape Corrig\'e d\'etaill\'e}
  \end{center}
  \vspace{0.2em}
  \hrule height 0.8pt
  \vspace{1.5pt}
  \hrule height 0.4pt
  \vspace{1.2em}}

% --- Encadres pedagogiques (corriges) --------------------------------
% Filet vertical gauche, texte legerement reduit, aucun fond : sobre,
% photocopiable, et tolere les sauts de page (package framed).
% Usage, avec parcimonie (2-3 encadres par corrige, la ou ils eclairent
% une demarche) :
%   \begin{methode} ... \end{methode}   -> « Méthode »
%   \begin{rappel}  ... \end{rappel}    -> « Rappel de cours »
%   \begin{piege}   ... \end{piege}     -> « Piège classique »
\usepackage{framed}
\newenvironment{pedagobarre}{%
  \def\FrameCommand{{\vrule width 1pt}\hspace{0.9em}}%
  \MakeFramed{\advance\hsize-\width\FrameRestore}}%
 {\endMakeFramed}
\newenvironment{blocpedago}[1]{%
  \par\medskip
  \begin{pedagobarre}\small
  \noindent{\bfseries\scshape #1.}\hspace{0.5em}\ignorespaces
}{%
  \end{pedagobarre}\medskip}
\newenvironment{methode}{\begin{blocpedago}{M\'ethode}}{\end{blocpedago}}
\newenvironment{rappel}{\begin{blocpedago}{Rappel de cours}}{\end{blocpedago}}
\newenvironment{piege}{\begin{blocpedago}{Pi\`ege classique}}{\end{blocpedago}}


% Macros maths / physique
\newcommand{\vv}[1]{\overrightarrow{#1}}
\newcommand{\R}{\mathbb{R}}
\newcommand{\N}{\mathbb{N}}
\newcommand{\vi}{\vec{\imath}}
\newcommand{\vj}{\vec{\jmath}}
% Chimie : formules en romain
\newcommand{\ch}[1]{\ensuremath{\mathrm{#1}}}

% Titre de sujet
\newcommand{\entetesujet}[1]{\begin{center}{\Large\bfseries #1}\end{center}\vspace{0.3em}\hrule\vspace{1em}}

% Bandeau de matiere : CHIMIE / PHYSIQUE avec bareme
\newcommand{\matiere}[2]{\par\bigskip\begin{center}\rule{2.2cm}{0.4pt}\quad{\large\bfseries #1}\ \textit{#2}\quad\rule{2.2cm}{0.4pt}\end{center}\nopagebreak\medskip}

% Titres d'exercices et parties
\newcommand{\exo}[2]{\medskip\noindent{\large\bfseries Exercice #1}\hfill\textit{#2}\par\nopagebreak\smallskip}
\newcommand{\partie}[1]{\medskip\noindent{\bfseries Partie #1}\par\nopagebreak\smallskip}
% Titre de question (corriges) : en gras, sur sa propre ligne
\newcommand{\question}[1]{\par\medskip\noindent{\bfseries #1}\par\nopagebreak\smallskip}

\setlist{leftmargin=2em,itemsep=2pt,topsep=3pt}


\begin{document}

\entetecorrige{2019}{C}{Physique-Chimie}

\matiere{CHIMIE}{8 points}
\partie{A : vérification des connaissances}

\question{1. Appariement}
\textbf{$1 = c$.} La constante d'acidité $K_a$ du couple $\ch{AH/A^-}$ ($\ch{AH + H_2O \rightleftharpoons A^- + H_3O^+}$) : $K_a = \dfrac{[\ch{A^-}][\ch{H_3O^+}]}{[\ch{AH}]}$.

\textbf{$2 = d$.} Loi de décroissance radioactive $N(t) = N_0 e^{-\lambda t}$. La demi-vie $T$ vérifie $N(t+T) = \frac{N(t)}{2}$, d'où $T = \dfrac{\ln 2}{\lambda}$ (soit $\lambda = \dfrac{\ln 2}{T}$).
\begin{center}
\begin{tikzpicture}[>=Stealth,scale=1,every node/.style={font=\footnotesize}]
  \draw[->] (0,0)--(5,0) node[right]{$t$}; \draw[->] (0,0)--(0,3.2) node[above]{$N(t)$};
  \draw[thick] plot[domain=0:4.8,samples=60] (\x,{3*exp(-0.55*\x)});
  \draw[dashed] (0,3)--(0,3); \node[left] at (0,3){$N_0$};
  \draw[dashed] (1,0)--(1,{3*exp(-0.55)})--(0,{3*exp(-0.55)}); \node[left] at (0,1.73){$N(t)$}; \node[below] at (1,0){$t$};
  \draw[dashed] (2.26,0)--(2.26,{3*exp(-0.55*2.26)})--(0,{3*exp(-0.55*2.26)}); \node[left] at (0,0.87){$\frac{N(t)}{2}$}; \node[below] at (2.26,0){$t{+}T$};
\end{tikzpicture}
\end{center}

\textbf{$3 = b$.} Pour $\alpha A + \beta B \rightleftharpoons \gamma C + \delta D$, la constante d'équilibre est $K = \dfrac{[C]_{eq}^{\gamma}[D]_{eq}^{\delta}}{[A]_{eq}^{\alpha}[B]_{eq}^{\beta}}$.

\textbf{$d = 1$.} Niveaux d'énergie de l'atome d'hydrogène : un électron absorbe un photon pour monter d'un niveau $p$ à $n > p$, ou en émet un pour descendre. $\Delta E = E_0\left(\dfrac{1}{p^2} - \dfrac{1}{n^2}\right) = \dfrac{hc}{\lambda}$, d'où $\dfrac{1}{\lambda} = \dfrac{E_0}{hc}\left(\dfrac{1}{p^2} - \dfrac{1}{n^2}\right)$ ; comparé à Balmer $\frac{1}{\lambda} = R_H\left(\frac{1}{p^2} - \frac{1}{n^2}\right)$, on a $R_H = \dfrac{E_0}{hc}$.
\begin{center}
\begin{tikzpicture}[>=Stealth,scale=0.95,every node/.style={font=\scriptsize}]
  \foreach \n/\y in {1/0,2/2.55,3/3.22,4/3.51,5/3.66,6/3.75}{\draw (0,\y)--(4.5,\y); \node[right] at (4.5,\y){$n=\n$};}
  \draw (0,3.9)--(4.5,3.9); \node[right] at (4.5,3.9){$n=\infty$};
  \node[left] at (0,0){$-13{,}6$}; \node[left] at (0,3.9){$0$}; \node[above left] at (0,3.9){$E$ (eV)};
  % Lyman -> n=1
  \foreach \yy in {2.55,3.22,3.51}{\draw[->] (0.5,\yy)--(0.5,0.05);}
  \node at (0.9,-0.35){Lyman (UV)};
  % Balmer -> n=2
  \foreach \yy in {3.22,3.51,3.66}{\draw[->] (2.2,\yy)--(2.2,2.6);}
  \node at (2.4,2.15){Balmer (visible)};
  % Paschen -> n=3
  \foreach \yy in {3.51,3.66,3.75}{\draw[->] (3.6,\yy)--(3.6,3.27);}
  \node at (3.5,3.05){Paschen (IR)};
\end{tikzpicture}
\end{center}

\question{2. Schéma à compléter}
$\ch{(C_2H_5)_2NH + H_2O \rightleftharpoons (C_2H_5)_2NH_2^+ + OH^-}$. L'atome d'azote porte un doublet non liant, ce qui confère aux amines un caractère basique (Br\o nsted-Lowry).

\question{3. Questions à choix multiples}
\textbf{$3.1 = a\cdot 3$.} Pour un mélange équimolaire, le rendement d'estérification (5\,\%, 60\,\% ou 67\,\%) diminue quand la classe de l'alcool augmente : pour un alcool \emph{secondaire}, il vaut \textbf{60\,\%}.

\textbf{$3.2 = b\cdot 3$.} $\lambda = \dfrac{hc}{E_n - E_p}$ est maximale si $E_n - E_p$ est minimale, c'est-à-dire pour $n = p+1$. Les plus petits écarts sont $E_2 - E_1$, $E_3 - E_2$, $E_4 - E_3$ ; le plus petit est $E_4 - E_3$. La plus grande longueur d'onde émise appartient donc à la \textbf{série de Paschen}.

\partie{B : application des connaissances}

\question{1.} $-\log C_0 = -\log(5\cdot10^{-2}) = 1{,}30$. Comme $\text{pH} > -\log C_0$, l'acide monochloroéthanoïque est un \textbf{acide faible}.

\question{2.} $\ch{CH_2ClCOOH + H_2O \rightleftharpoons CH_2ClCOO^- + H_3O^+}$.

\question{3. a.} Autoprotolyse (acide faible, dissociation limitée) :
\[ \ch{CH_2ClCOOH + H_2O \rightleftharpoons CH_2ClCOO^- + H_3O^+} \;;\quad \ch{2\,H_2O \rightleftharpoons H_3O^+ + HO^-} \]
\underline{Espèces :} ions \ch{H_3O^+}, \ch{OH^-}, \ch{CH_2ClCOO^-} ; molécules \ch{CH_2ClCOOH}, \ch{H_2O}.

$\bullet$ $[\ch{H_2O}] = \dfrac{1000}{18} = 55{,}6\ \text{mol}\cdot\text{L}^{-1}$. \quad $\bullet$ $[\ch{H_3O^+}] = 10^{-2{,}1} = 7{,}94\cdot10^{-3}\ \text{mol}\cdot\text{L}^{-1}$.

$\bullet$ $[\ch{OH^-}] = \dfrac{K_e}{[\ch{H_3O^+}]} = \dfrac{10^{-14}}{7{,}94\cdot10^{-3}} = 1{,}26\cdot10^{-12}\ \text{mol}\cdot\text{L}^{-1}$.

$\bullet$ Électroneutralité $[\ch{H_3O^+}] = [\ch{CH_2ClCOO^-}] + [\ch{OH^-}]$, avec $[\ch{OH^-}] \ll [\ch{H_3O^+}]$ : $[\ch{CH_2ClCOO^-}] = 7{,}94\cdot10^{-3}\ \text{mol}\cdot\text{L}^{-1}$.

$\bullet$ Conservation $C_A = [\ch{CH_2ClCOOH}] + [\ch{CH_2ClCOO^-}]$ : $[\ch{CH_2ClCOOH}] = 5\cdot10^{-2} - 7{,}94\cdot10^{-3} = 4{,}21\cdot10^{-2}\ \text{mol}\cdot\text{L}^{-1}$.

\textbf{b.} $\text{p}K_a = \text{pH} - \log\dfrac{[\ch{CH_2ClCOO^-}]}{[\ch{CH_2ClCOOH}]} = 2{,}1 - \log\dfrac{7{,}94\cdot10^{-3}}{4{,}21\cdot10^{-2}} = 2{,}82$.

\question{4.} Dosage : $\ch{CH_2ClCOOH + OH^- \longrightarrow CH_2ClCOO^- + H_2O}$. À la demi-équivalence $\text{pH} = \text{p}K_a$ et $V_B = \dfrac{V_{BE}}{2}$. À l'équivalence $C_A V_A = C_B V_{BE}$, d'où $V_B = \dfrac{C_A V_A}{2C_B} = \dfrac{5\cdot10^{-2}\times 20\cdot10^{-3}}{2\times 0{,}1} = 5\cdot10^{-3}$ L $= 5$ mL.

\matiere{PHYSIQUE}{12 points}
\partie{A : vérification des connaissances}

\question{1. Questions à choix multiples}
\textbf{a. Vrai.} Pendule conique : boule de masse $m$, fil de longueur $l$, en mouvement circulaire uniforme autour de $(\Delta)$.
\begin{center}
\begin{tikzpicture}[>=Stealth,scale=1,every node/.style={font=\footnotesize}]
  \draw[dashed] (0,3)--(0,-0.4); \node[left] at (0,-0.2){$(\Delta)$};
  \coordinate (S) at (0,3); \fill (S) circle (1pt);
  \coordinate (M) at (1.6,0.6);
  \draw[thick] (S)--(M); \fill (M) circle (2pt);
  \draw[dashed] (M)++(-1.6,0) ellipse (1.6 and 0.4);
  \draw (0,2.3) arc[start angle=-90,end angle=-25,radius=0.7]; \node at (0.35,2.05){\scriptsize$\theta$};
  \node[left] at (0.9,1.7){$l$};
  \draw[->,thick,blue] (M)--($(S)!0.55!(M)$) node[midway,above right]{$\vv{T}$};
  \draw[->,thick,green!60!black] (M)--++(0,-1) node[below]{$\vv{P}$};
  \draw[->,thick,red] (M)--++(-1,0) node[left]{$m\vv{a}$};
  \draw[dashed] (M)++(-1.6,0)--(M); \node[below] at (0.7,0.55){\scriptsize$R$};
\end{tikzpicture}
\end{center}
$a_T = 0$ (uniforme), $a = a_N = R\omega^2$ avec $R = l\sin\theta$, donc $a = l\omega^2\sin\theta$. PFD : $\vv{P} + \vv{T} = m\vv{a}$ donne $\begin{cases} T\sin\theta = ml\omega^2\sin\theta \\ T\cos\theta = mg \end{cases}$ ; en divisant : $\dfrac{1}{\cos\theta} = \dfrac{l\omega^2}{g}$.

\textbf{b. Vrai.} L'effet photoélectrique est l'émission d'électrons par un métal recevant de l'énergie rayonnante.

\textbf{c. Vrai.} En mouvement circulaire uniforme, l'accélération est normale et vaut $R\omega^2$ ($a_T = 0$).

\textbf{d. Faux.} Un $RLC$ série est en résonance quand $X_L = X_C$ ; alors $Z = \sqrt{R^2 + (X_L - X_C)^2}$ est minimale et vaut $R$.

\question{2. Texte à trous}
Un ébranlement est \emph{transversal} lorsque la perturbation du milieu élastique est \emph{perpendiculaire} à la direction de propagation.

\partie{B : application des connaissances}

\question{1.} Courbe 1 ($y$--$t$) : période $T = 2\cdot10^{-3}$ s. Courbe 2 ($y$--$x$) : longueur d'onde $\lambda$. D'où $v = \dfrac{\lambda}{T} = \dfrac{\lambda}{2\cdot10^{-3}} = 500\lambda\ (\text{m}\cdot\text{s}^{-1})$.
\begin{center}
\begin{tikzpicture}[>=Stealth,scale=0.8,every node/.style={font=\scriptsize}]
  \begin{scope}
  \draw[->] (0,0)--(5,0) node[right]{$t$ (s)}; \draw[->] (0,-1.3)--(0,1.5) node[above]{$y$ (mm)};
  \draw[dashed] (0,1)--(4,1); \draw[dashed] (0,-1)--(4,-1);
  \node[left] at (0,1){$0{,}2$}; \node[left] at (0,-1){$-0{,}2$};
  \draw[thick] plot[domain=0:4,samples=80] (\x,{sin(deg(pi*\x))});
  \draw[<->] (0,1.25)--(2,1.25); \node[above] at (1,1.25){$T$};
  \node at (2.2,-1.7){Courbe 1};
  \end{scope}
  \begin{scope}[xshift=6.2cm]
  \draw[->] (0,0)--(5,0) node[right]{$x$}; \draw[->] (0,-1.3)--(0,1.5) node[above]{$y$ (mm)};
  \draw[dashed] (0,1)--(4,1); \draw[dashed] (0,-1)--(4,-1);
  \node[left] at (0,1){$0{,}2$}; \node[left] at (0,-1){$-0{,}2$};
  \draw[thick] plot[domain=0:4,samples=80] (\x,{sin(deg(pi*\x))});
  \draw[<->] (0,1.25)--(2,1.25); \node[above] at (1,1.25){$\lambda$};
  \draw[<->] (0,-1.25)--(3,-1.25); \node[below] at (1.5,-1.25){$d_1$};
  \node at (2.2,-1.7){Courbe 2};
  \end{scope}
\end{tikzpicture}
\end{center}

\question{2.} $y_O(t) = a_m\sin(\omega t + \varphi)$ avec $\omega = \dfrac{2\pi}{T} = 1000\pi$, $a = 2\cdot10^{-4}$ m. Comme $y_O(0) = 0$, $\varphi = 0$ : $y_O(t) = 2\cdot10^{-4}\sin(1000\pi t)$.

\question{3.} $M$ reproduit $O$ avec le retard $t_M = \dfrac{d}{v} = \dfrac{12{,}5}{500\lambda}$ : $y_M(t) = y_O(t - t_M) = 2\cdot10^{-4}\sin\left(1000\pi t - \dfrac{25\pi}{\lambda}\right)$.

\question{4.} $\dfrac{\lambda}{T} = \dfrac{d_1}{t_1}$, d'où $t_1 = \dfrac{d_1}{\lambda}T$. La courbe 2 donne $\dfrac{d_1}{\lambda} = \dfrac{6}{4}$, donc $t_1 = \dfrac{6}{4}T = 3\cdot10^{-3}$ s.

\partie{C : résolution d'un problème}
\begin{center}
\begin{tikzpicture}[>=Stealth,scale=1.1,every node/.style={font=\footnotesize}]
  \coordinate (O) at (0,2.2); \coordinate (B) at (0,0); \coordinate (Ap) at (2.2,2.2);
  % arc AB (quart de cercle centre O rayon 2.2)
  \draw[thick] (B) arc[start angle=-90,end angle=0,radius=2.2];
  \draw[thick] (B)--(-1.8,0); \node[below] at (-1.8,0){$C$};
  \draw[dashed] (O)--(B); \draw[dashed] (O)--(Ap);
  \coordinate (A) at (2.2,2.2); \fill (A) circle (1.2pt) node[right]{$A$};
  \fill (O) circle (1.2pt) node[above]{$O$}; \fill (B) circle (1.2pt) node[below]{$B$};
  \draw (0,1.5) arc[start angle=-90,end angle=0,radius=0.7]; \node at (0.32,1.35){\scriptsize$\theta$};
  \coordinate (I) at (0,2.2); \draw[dashed] (2.2,2.2)--(0,2.2);
  \fill (0,1.55) circle (1pt) node[left]{$I$}; \draw[dashed] (2.2,1.55)--(0,1.55);
  \draw[<->] (-0.18,0)--(-0.18,1.55); \node[left] at (-0.18,0.78){$h$};
  % forces en A
  \draw[->,thick,green!60!black] (A)--++(0,-1) node[below]{$\vv{P}$};
  \draw[->,thick,red] (A)--++(-0.8,0.8) node[above left]{$\vv{R}$};
  \draw[->,thick,red!60!black] (A)--++(-0.9,0.55) node[left]{$\vv{R_N}$};
  \draw[->,thick] (A)--++(0.15,0.6) node[above]{$\vv{f}$};
\end{tikzpicture}
\end{center}

\question{1. a.} Référentiel terrestre galiléen, système = solide de masse $m$. TEC entre $A$ et $B$ : $\frac{1}{2}mv_B^2 - \frac{1}{2}mv_A^2 = W_{\vv{P}} + W_{\vv{R}}$. La réaction normale ne travaille pas ($\vv{R_N} \perp \vv{dl}$), $W_{\vv{f}} = -f\cdot\widehat{AB}$. Avec $v_A = 0$, $\widehat{AB} = r\theta$, $OI = r\cos\theta$ et $h = r(1-\cos\theta)$ :
\[ \frac{1}{2}mv_B^2 = mgr(1-\cos\theta) - f r\theta \;\Rightarrow\; v_B = \sqrt{2r\left[g(1-\cos\theta) - \frac{f}{m}\theta\right]}. \]

\textbf{b.} Entre $B$ et $C$ (horizontal, $W_{\vv{P}} = 0$) : $\frac{1}{2}mv_C^2 - \frac{1}{2}mv_B^2 = -f\cdot BC = -2fr$, d'où
\[ v_C = \sqrt{2r\left[g(1-\cos\theta) - \frac{f}{m}(\theta + 2)\right]}. \]

\question{2. a.} En $C$, $v_C = 0$ : $g(1-\cos\theta) - \dfrac{f}{m}(\theta + 2) = 0$, d'où $f = \dfrac{mg(1-\cos\theta)}{\theta + 2}$.

\textbf{b.} $f = \dfrac{60\times 10\times(1 - \cos\frac{\pi}{4})}{\frac{\pi}{4} + 2} \approx 63{,}1$ N.

\textbf{c.} $v_B \approx \sqrt{2\times 5\left[10\left(1 - \cos\frac{\pi}{4}\right) - \frac{63{,}1}{60}\times\frac{\pi}{4}\right]} \approx 4{,}59\ \text{m}\cdot\text{s}^{-1}$.

\end{document}
